% p.36 — الباب التاسع العشر في اصطلاحات التعليل
% Layout (as in the lithograph): under the chapter title, one `tables` row with two tables
% side by side. RIGHT: terminology (term in the right column, with the example letter/word
% نَصَرَ، ن، ص، ر in the same column for the first four rows; definition in the left column).
% LEFT: a three-column table وزن / اصل / بدل. Below both, a paragraph of Arabic prose
% spanning the page.
\mspage{36}
\bab{الباب التاسع العشر في اصطلاحات التعليل}
\begin{tables}
\begin{sarf}{|c|p{58mm}|}
فَعَلَ \quad نَصَرَ & الثلاثي المجرد ـ الفعل الماضي\\
الفاء \quad ن & الحرف الاول في ماضي الثلاثي المجرد مثل النون في نَصَرَ\\
العين \quad ص & الحرف الثاني في ماضي الثلاثي المجرد مثل الصاد في نَصَرَ\\
اللام \quad ر & الحرف الثالث في ماضي الثلاثي المجرد مثل الراء في نَصَرَ\\
حروف العلة & وَاوٌ اَلِفٌ يَاءٌ\\
مَهْمُوْز & ما فيه هَمْزَةٌ اَخَذَ سَأَلَ قَرَأَ\\
مَعْلُوْل & ما فيه حرف العلة وَعَدَ يَسَرَ قَالَ\\
مُضَاعَف & ما فيه تشديد مَلَّ فَرَّ عَضَّ\\
مُضَاعَف & ما تَكَرَّرَ الفاء والعين مثل زَلْزَلَ\\
صَحِيْح & ما ليس فيه همزة ولا تشديد ولا حرف العلة مثل نَصَرَ\\
مِثَال & ما في الفاء حرف علة مثل وَعَدَ يَسَرَ\\
اجوف & ما في العين حرف علة مثل قَوَلَ بَيَعَ\\
نَاقِص & ما في اللام حرف علة مثل دَعَوَ رَمَيَ\\
لفيف مفروق & ما في الفاء واللام حرف علة مثل وَقَيَ\\
لفيف مقرون & ما في العين واللام حرف علة مثل عَوَيَ\\
\end{sarf}
\begin{sarf}{|c|c|c|}
\thead{وزن} & \thead{اصل} & \thead{بدل}\\\hline
نَصَرَ & قَوَلَ & قَالَ\\
نَصَرْتُ & قَوَلْتُ & قُلْتُ\\
اُنْصُرْ & اُقْوُلْ & قُلْ\\
ضَرَبَ & بَيَعَ & بَاعَ\\
مَضْرُوْبٌ & مَبْيُوْعٌ & مَبِيْعٌ\\
لَمْ يَضْرِبْنَ & لَمْ يَبْيِعْنَ & لَمْ يَبِعْنَ\\
اَفْعَلَ & اَعْوَلَ & اَعَانَ\\
لَيُفْعِلَنَّ & لَيُعْوِنَنَّ & لَيُعِيْنَنَّ\\
مُفْعِلُوْنَ & مُعْوِنُوْنَ & مُعِيْنُوْنَ\\
اُنْصُرْ & اُمْدُدْ & مُدَّ\\
نَاصِرُوْنَ & مَادِدُوْنَ & مَادُّوْنَ\\
ضَرَبْتُ & بَيَعْتُ & بِعْتُ\\
لَمْ يَنْصُرْ & لَمْ يَقْوُلْ & لَمْ يَقُلْ\\
لَا اُفْعِلَنَّ & لَا اُعْوِنَنَّ & لَا اُعِيْنَنَّ\\
\end{sarf}
\end{tables}
\pstart
نكتب في ما يأتي جميع الافعال المعتلة المشهورة ونشير الى ابوابها واوزانها واصولها وتغيراتها يلزم على الاستاذ اولا ان يمرن الطلبة في هذه الافعال حتى يكون السنتهم مُذَلَّلَةً في التصريف واذا \com{اعتاد والتصريف}{كذا في الأصل، ولعل الصواب: اعتادوا التصريف} فيعلم لهم ثانيا الابواب والاوزان والاصول والمُبَدَّل كما ذكر قبل تفصيلا من الصرف الصغير والكبير وغيرها حتى يكون لهم ملكة في التمييز بين الاصل والباب والوزن والتبديل مثلا اَجَابَ اصله اَجْوَبَ وبابه اَفْعَلَ يُفْعِلُ والمبدل اَجَابَ كذالك قُلْتُ اصله قَوَلْتُ وبابه نَصَرَ يَنْصُرُ ووزنه فَعَلْتُ والمبدل قُلْتُ كذالك نَمُدُّ اصله نَمْدُدُ ووزنه نَفْعُلُ وبابه نَصَرَ يَنْصُرُ كذالك تقيس في سائر التصريفات المذكورة في الجدول كما مثلنا في الفوق يسيرا
\pend
